\section{cMCP Design}
\label{sec:design}

\subsection{Architecture Overview}

cMCP is a transparent MCP proxy. An MCP client connects to cMCP as if connecting to an MCP
server; cMCP connects to one or more upstream MCP servers as an MCP client. The client
requires no modification. Figure~\ref{fig:arch} sketches the component layout.

\begin{figure}[h]
\centering
\begin{tikzpicture}[
  font=\footnotesize,
  >={Stealth[round]},
  endpoint/.style={draw, rounded corners, minimum width=2.2cm, minimum height=2.6em, align=center, fill=gray!8},
  stage/.style={draw, rounded corners, minimum width=4.0cm, minimum height=2.0em, align=center, fill=blue!5},
  tee/.style={draw, dashed, thick, rounded corners, draw=blue!55!black, inner sep=12pt},
  note/.style={align=center, font=\scriptsize\itshape, text=black!70},
]
% Gateway internal pipeline (vertical)
\node[stage] (catalog) {1.\ Catalog check (hash-pinned)};
\node[stage, below=0.45cm of catalog] (cedar) {2.\ Cedar policy evaluation};
\node[stage, below=0.45cm of cedar] (dispatch) {3.\ Dispatch};
\node[stage, below=0.45cm of dispatch] (inspect) {4, 5.\ Inspect + session sensitivity};
\node[stage, below=0.45cm of inspect] (audit) {6, 7.\ Audit chain + TRACE Claim};

% TEE boundary
\begin{scope}[on background layer]
  \node[tee, fit=(catalog)(cedar)(dispatch)(inspect)(audit),
        label={[text=blue!55!black]above left:\textbf{cMCP Gateway (TEE)}}] (gw) {};
\end{scope}

% External endpoints
\node[endpoint, left=1.4cm of cedar] (client) {MCP\\Client};
\node[endpoint, right=1.4cm of dispatch] (server) {MCP\\Server(s)};
\node[endpoint, below=1.0cm of audit, fill=green!6] (verifier) {Verifier};

% Request / response flow
\draw[->] (client.east) -- (catalog.west);
\draw[->] (catalog) -- (cedar);
\draw[->] (cedar) -- (dispatch);
\draw[->] (dispatch.east) -- node[note, above, pos=0.75]{call} (server.west);
\draw[->] (server.south) to[out=270,in=0] node[note, right, pos=0.3]{response} (inspect.east);
\draw[->] (inspect) -- (audit);

% Startup measurement and output
\node[note, above=0.9cm of gw.north] (meas)
  {bound at startup: key thumbprint $\|$ salt $\rightarrow$ \texttt{REPORT\_DATA}; per session: thumbprint $\|$ H(chain root)};
\draw[->, dotted, thick] (meas) -- (gw.north);
\draw[->] (audit.south) -- node[note, right, pos=0.75]{signed TRACE Claim} (verifier.north);
\end{tikzpicture}
\caption{cMCP interposes on every tool call between an MCP client and upstream MCP
servers. The numbered stages correspond to the tool-call flow in
Section~\ref{sec:design}. The gateway key is bound into the attestation report at startup,
and each session's report also commits the audit-chain root; policy-bundle and catalog
hashes travel in the signed TRACE Claim, which lets an external verifier confirm what
governed the session.}
\label{fig:arch}
\end{figure}

At startup, cMCP performs four initialization steps before accepting any connections: (1) it
generates a fresh Ed25519 keypair whose private half never leaves TEE memory; (2) it loads
the Cedar policy bundle, records its SHA-256 hash, and refuses to start if the hash differs
from a configured expected value; (3) it loads the operator-approved tool catalog file,
records the aggregate catalog hash, and applies the same expected-hash check; and (4) it
obtains an attestation report from the TEE with \texttt{REPORT\_DATA} set to the key
nonce described below. The policy and catalog hashes are carried in
the signed claim body rather than in \texttt{REPORT\_DATA} at the measured revision.

\subsection{Tool Call Flow}

When the MCP client sends a \texttt{tools/call} request, cMCP executes the following
sequence:

\begin{enumerate}[leftmargin=1.5em]
  \item \textbf{Catalog check.} Verify that the requested tool is in the pinned catalog and
    that the in-memory catalog hash still equals the startup hash. If it does not, the
    session is marked \texttt{catalog\_drift} and the call is refused.
  \item \textbf{Cedar evaluation.} Evaluate the Cedar policy bundle against the call
    context. If the decision is \texttt{forbid}, return a policy-denial response immediately
    without forwarding the call.
  \item \textbf{Dispatch.} Forward the call to the upstream MCP server.
  \item \textbf{Inspection.} Pass the response through the \texttt{InspectionPipeline},
    which classifies sensitivity tags (e.g., \texttt{hipaa\_phi}, \texttt{pii}) and detects
    prompt-injection patterns.
  \item \textbf{Session state update.} Call \texttt{update\_from\_inspection()} on the
    session's \texttt{SessionState}, which may raise the session sensitivity level.
  \item \textbf{Audit append.} Append a hash-chained entry to the session's audit log,
    recording the tool name, policy decision, sensitivity tags, and the previous entry's
    hash.
  \item \textbf{Call log record.} Record the call in the session's \texttt{SessionCallLog}
    with a monotonically increasing sequence number, for temporal adjacency tracking.
\end{enumerate}

Steps 1 and 2 must both pass for the call to proceed. A Cedar allow decision does not
bypass the catalog check, and a clean catalog does not bypass Cedar.

\subsection{Session Sensitivity State Machine}

Each session maintains a \texttt{SessionState} object with a \texttt{max\_sensitivity}
field. The sensitivity order is:

\[
  \text{public}(0) < \text{pii}(1) < \text{confidential}(2) < \text{hipaa\_phi} = \text{mnpi} = \text{trade\_secret}(3)
\]

\texttt{max\_sensitivity} is initialized to \texttt{public} and transitions only upward:
whenever the inspection pipeline tags a response with a higher sensitivity level, the state
machine raises \texttt{max\_sensitivity} to that level and records which call triggered the
transition. The session sensitivity never decreases automatically. An operator reset, which
rotates the session identifier, is the only way to lower it.

This monotonic design is the source of both its guarantee and its cost. It guarantees that
once PHI enters the session context, every subsequent external call is attributed as
potentially PHI-adjacent. It causes false positives when a subsequent external call is
provably unrelated to the PHI. Calls can now declare a sensitivity class through
\texttt{\_cmcp.data\_class}, but the declaration is agent-cooperative input rather than an
independently proven lineage fact.

\subsection{Tool Catalog Pinning}

At startup, cMCP computes a definition hash for each tool:

\[
  h_{\text{tool}} = \text{SHA-256}(\text{tool\_name} \| \text{description} \| \text{inputSchema})
\]

where the fields are serialized as canonical JSON with sorted keys. The aggregate catalog
hash is SHA-256 of the sorted concatenation of all per-tool hashes. Both hashes are stored
in TEE-protected memory.

The approved catalog is an operator-curated file, not a live fetch. A catalog whose hash
differs from the expected value raises \texttt{CatalogHashMismatch} at load time and the
gateway does not start. Clients receive \texttt{tools/list} from the approved catalog, so a
description an upstream server changes after approval never reaches the agent. Before every
\texttt{tools/call}, the gateway compares the in-memory catalog hash with the startup value
and refuses the call on a mismatch. At the measured revision it does not compare the
upstream server's own \texttt{tools/list} with the approved entry; that comparison was added
later (Section~\ref{sec:eval}, Limitations). Hashing a definition cannot detect a server
that keeps its advertised definition and changes its behavior.

\subsection{TRACE Claim Generation}

At session close (or on demand), cMCP produces a TRACE Claim by:

\begin{enumerate}[leftmargin=1.5em]
  \item Computing the session nonce:
    $\text{nonce} = \text{TP}(\text{pubkey}) \,\|\, \text{SHA-256}(\text{chain\_root})$, the
    32-byte RFC~7638 JWK Thumbprint of the gateway public key followed by the hash of the
    session's first audit entry, which records the session identifier. The startup report
    uses a 32-byte random salt in place of the chain-root hash.
  \item Requesting an attestation report from the TEE with \texttt{REPORT\_DATA} set to the
    64-byte nonce (which fills the field exactly). This happens when the session opens; if
    the call fails, the claim falls back to the startup report and the verifier reports the
    chain-root binding as missing.
  \item Assembling the \texttt{RuntimeClaim} Pydantic model with the attestation report,
    policy bundle hash, catalog hash, session call summary, and audit chain tip.
  \item Signing the canonical JSON serialization (sorted keys, no whitespace) of the claim
    body with the TEE's Ed25519 private key.
  \item Embedding the signature as a base64url-encoded field in the claim.
\end{enumerate}

The nonce construction binds the attestation report to the TEE instance: the first 32
bytes are the RFC~7638 JWK Thumbprint of the gateway public key, which a verifier
re-derives from the claim's \texttt{cnf.jwk} and compares to \texttt{REPORT\_DATA}$[{:}32]$,
confirming the report belongs to the key that signed the claim. The second half,
\texttt{REPORT\_DATA}$[32{:}64]$, commits the session's audit-chain root, so in hardware
mode the session is bound twice: by the hardware report and by the Ed25519 signature over
\texttt{session\_id}. Without the chain-root commitment an operator holding the signing
path could rebuild an internally consistent chain and re-sign it; this binding was added
in cMCP PR~\#368, before the measured revision.

\subsection{Cross-Organizational Attestation (Phase 2)}

When an MCP server also runs inside a TEE, the dual-TEE protocol extends the trust chain to
the server side (Figure~\ref{fig:dual-tee}). In Phase~2:

\begin{itemize}[leftmargin=1.5em]
  \item The server TEE generates its own Ed25519 keypair (Phase~2 key).
  \item At session start, both TEEs exchange public keys and agree on the shared session
    identifier.
  \item The Phase~2 nonce is the \emph{server} key thumbprint followed by its own salt,
    structurally identical to the Phase~1 nonce but bound to a different key. Both phases
    carry the shared session identifier in their signed claim bodies.
  \item The server TEE produces an independent TRACE Claim signed with the Phase~2 key.
\end{itemize}

\begin{figure}[h]
\centering
\begin{tikzpicture}[
  font=\footnotesize,
  >={Stealth[round]},
  tee/.style={draw, dashed, thick, rounded corners, align=center, minimum width=4.3cm,
              minimum height=4.2em, inner sep=8pt},
  claim/.style={draw, rounded corners, fill=blue!5, align=center, minimum height=2.0em},
  note/.style={font=\scriptsize\itshape, text=black!70},
]
\node[tee, draw=blue!55!black] (gw)
  {\textbf{Phase 1: Gateway TEE}\\[2pt] key $K_1$\\ $n_1=\text{TP}(K_{1,\text{pub}})\,\|\,s_1$};
\node[tee, draw=orange!75!black, right=3.0cm of gw] (srv)
  {\textbf{Phase 2: Server TEE}\\[2pt] key $K_2$\\ $n_2=\text{TP}(K_{2,\text{pub}})\,\|\,s_2$};

% shared session identifier link
\draw[<->, thick] (gw.north) to[out=60,in=120]
  node[note, above, pos=0.5]{shared session id $\sigma$} (srv.north);

\node[claim, below=0.9cm of gw] (c1) {TRACE Claim 1\\ signed by $K_1$};
\node[claim, below=0.9cm of srv] (c2) {TRACE Claim 2\\ signed by $K_2$};
\draw[->] (gw) -- (c1);
\draw[->] (srv) -- (c2);

\node[draw, rounded corners, fill=green!6, align=center, minimum height=2.0em,
      below=0.9cm of $(c1)!0.5!(c2)$] (ver) {Verifier:\ links via $\sigma$,\ verifies $K_1,K_2$ independently};
\draw[->] (c1.south) -- (ver.north west);
\draw[->] (c2.south) -- (ver.north east);
\end{tikzpicture}
\caption{The dual-TEE attestation chain (Property~\ref{prop:cross-org}). Two
independently generated TEE keys produce two TRACE Claims linked by a shared
session identifier $\sigma$ but verifiable independently: tampering with one
claim does not affect the other.}
\label{fig:dual-tee}
\end{figure}

A verifier who holds both claims can confirm: (a) both claims reference the same session
identifier; (b) the Phase~1 nonce is bound to the gateway TEE key; (c) the Phase~2 nonce
is bound to the server TEE key; (d) tampering with one claim does not affect the validity
of the other. The two claims are linked but independently verifiable.

Phase~2 server TEE deployment is in progress; the experiment in Section~\ref{sec:eval}
simulates the protocol in software.
